翻开这一卷，我们就要带着前面练过的基本功，开始解决视觉组里的具体问题了。

卷一里，我们从装好 Linux、熟悉命令行开始，学着用 Git 管理修改，用现代 C++ 组织数据、处理资源与并发，再用 CMake 把几个源文件构建成一个能运行的工程。走到这里，你应该已经能在自己的电脑上读懂、修改、构建一份程序，也能把修改和运行结果交给队友。如果这些事情还做不顺，先回去补一补，完成卷一的通识验收。后面的工具和算法都会用到它们，基础上的疑问不会因为翻过一页就自动消失。

不过，工程能编译，离解决实际问题还有一段路。相机拍到了画面，怎样找出我们关心的目标？有了目标的位置，怎样处理测量的抖动，又怎样把结果及时交给下一环？机器人怎样判断自己在哪、该怎么走，又怎样根据比赛状态决定下一步做什么？本卷会从这些问题出发，带你学习 ROS2、图像处理、深度学习，以及通信、成像、滤波、几何与控制等专项技术。卷一让我们有能力动手，这一卷要让动手做出的东西有依据、经得起检查。

组内按自瞄、导航、决策和雷达四个方向开展工作，它们会互相配合，你也不需要一口气学完所有章节。先和学长确认当前的任务，再安排学习顺序：自瞄从 ROS2、OpenCV 和相机入手，逐步走向识别、滤波、几何与控制；导航从 ROS2 走向定位、建图与路径执行；决策可以从 ROS2 的任务接口和导航章的行为树入手，关注怎样结合比赛状态、自身状态与感知结果选择任务，以及收到完成、失败或超时反馈后该怎么办；雷达则关注全场目标识别、坐标转换、多机通信和 SDR 接收解析。OpenCV、深度学习和 CUDA 等工具会在不同方向里用到，按任务补齐即可，标 $*$ 的小节也可以等需要时再读。选方向是给学习找一个落脚点，今后遇到问题，仍然要愿意走出自己熟悉的那几章。

这一卷里的陌生名词和公式会多一些，第一次读不顺很正常。先弄清它要解决什么问题，输入是什么，输出是什么，再跟着示例看它怎样工作。公式里的一个假设，可能就决定了方法能不能用在眼前的数据上；代码里一个参数的单位，也可能决定结果为什么总差一截。没想明白时，多画一张图、多检查一个中间结果，比把整段代码抄进工程里更有用。

学的时候，请给动手和核对结果留出时间。带编号的示例会说明操作步骤与预期结果，正文里的短代码片段则用来解释接口；Python 实验按说明准备 Python 3 和相应依赖。先跑通示例，再改一个条件，看看结果是否符合你的判断，最后独立完成练习和章末验收。遇到问题，把输入、配置、现象和已经尝试过的方法留下来。向学长求助时有据可查，将来换一位队友接手时，也少一些“当时怎么就能跑”的谜团。

也请记住，示例跑通只是开始。离线数据上的结果，换到真实相机、通信链路和机器人上，还要重新核对。不要急着给一次成功下结论，也不要因为一次失败就把所有参数改一遍。先确认哪一步出了变化，再提出猜测、动手验证，硬件实验就在对应设备上完成验收。这样的过程有时慢，却能让你知道自己究竟解决了什么。

等这些模块逐渐熟悉起来，卷三会带我们把它们放回自瞄、导航、决策和雷达的比赛链路，看看感知结果怎样支持任务选择，任务怎样交给导航等模块执行，接口怎样约定，问题怎样定位，又怎样和机械、电控的队友一起把系统调好。我们希望你带过去的，除了几个能运行的程序，还有解释结果、承认疑问和把事情做完的耐心。接下来，就从手头的一个问题开始。
